% !TeX program = xelatex
\documentclass[12pt,a4paper]{article}
\usepackage[margin=2cm]{geometry}
\usepackage{fontspec}
\usepackage{amsmath,amssymb}
\usepackage{enumitem}
\setmainfont{Times New Roman}
\setlength{\parindent}{0pt}
\setlength{\parskip}{6pt}
\begin{document}
\begin{center}
{\Large\bfseries Đề thi tuyển sinh vào lớp 10 TP. Hà Nội năm học 2016--2017}\\
{\large Phần Đại số}
\end{center}

\textbf{Câu 1.} Cho hai biểu thức
\[
A=\frac{7}{\sqrt{x}+8},\qquad
B=\frac{\sqrt{x}}{\sqrt{x}-3}+\frac{2\sqrt{x}-24}{x-9},
\qquad x\ge 0,\ x\ne 9.
\]
\begin{enumerate}[label=\alph*)]
\item Tính giá trị của biểu thức $A$ khi $x=25$.
\item Chứng minh rằng $B=\dfrac{\sqrt{x}+8}{\sqrt{x}+3}$.
\item Tìm $x$ để biểu thức $P=A\cdot B$ có giá trị nguyên.
\end{enumerate}

\textbf{Câu 2.} Giải bài toán sau bằng cách lập phương trình hoặc hệ phương trình:

Một mảnh vườn hình chữ nhật có diện tích $720\,\mathrm{m}^2$. Nếu tăng chiều dài thêm $10\,\mathrm m$ và giảm chiều rộng $6\,\mathrm m$ thì diện tích mảnh vườn không đổi. Tính chiều dài và chiều rộng của mảnh vườn.

\textbf{Câu 3.}
\begin{enumerate}[label=\alph*)]
\item Giải hệ phương trình
\[
\begin{cases}
\dfrac{3x}{x-1}-\dfrac{2}{y+2}=4,\\[4pt]
\dfrac{2x}{x-1}+\dfrac{1}{y+2}=5.
\end{cases}
\]
\item Trong mặt phẳng tọa độ $Oxy$, cho đường thẳng $(d): y=3x+m^2-1$ và parabol $(P):y=x^2$.
\begin{enumerate}[label=(\alph*)]
\item Chứng minh rằng $(d)$ luôn cắt $(P)$ tại hai điểm phân biệt với mọi $m$.
\item Gọi $x_1,x_2$ là hoành độ giao điểm của $(d)$ và $(P)$. Tìm $m$ để $(x_1+1)(x_2+1)=1$.
\end{enumerate}
\end{enumerate}

\textbf{Câu 5.} Với các số thực $x,y$ thỏa mãn
\[
 x-\sqrt{x+6}=\sqrt{y+6}-y,
\]
tìm giá trị lớn nhất và giá trị nhỏ nhất của biểu thức $P=x+y$.

\end{document}
